\section{Linear Kalman Filter State Estimation}

The \texttt{KalmanFilterEstimate} class implements a discrete-time linear Kalman filter for estimating the floating-base position and velocity of a legged robot. The filter fuses information from:

\begin{itemize}
\item IMU linear acceleration,
\item Joint encoder measurements,
\item Foot kinematics,
\item Foot contact states.
\end{itemize}

\subsection{State Vector}

Assume the robot has $(n_c)$ contact points. The state vector is defined as

\begin{equation}
\mathbf{x} =\begin{bmatrix}
\mathbf{p}_b \
\mathbf{v}_b \
\mathbf{p}_{f_1} \
\hdots \
\mathbf{p}_{f_{n_c}}
\end{bmatrix}
\in \mathbb{R}^{6+3n_c},
\end{equation}

where

\begin{itemize}
\item ($\mathbf{p}_b \in \mathbb{R}^{3})$ is the base cartesian position,
\item ($\mathbf{v}_b \in \mathbb{R}^{3})$ is the base cartesian/linear velocity,
\item ($\mathbf{p}_{f_i} \in \mathbb{R}^{3})$ is the world-frame position of foot (i).
\end{itemize}

Therefore,

\begin{equation}
n_x = 6 + 3n_c .
\end{equation}

\subsection{System Dynamics}

The state propagation model is

\begin{equation}
\mathbf{x}_{k+1} = \mathbf{A}\mathbf{x}_k
+
\mathbf{B}\mathbf{a}_k
+
\mathbf{w}_{k},
\end{equation}

where

\begin{equation}
\mathbf{a}_k
= \mathbf{R}_{WB}\mathbf{a}_{imu}
+
\mathbf{g}
\end{equation}

is the IMU acceleration expressed in the world frame.

The state transition matrix is

\begin{equation}
\mathbf{A}=
\begin{bmatrix}
\mathbf{I}_3 & \Delta t \mathbf{I}_3 & \mathbf{0} \\
\mathbf{0} & \mathbf{I}_3 & \mathbf{0} \\
\mathbf{0} & \mathbf{0} & \mathbf{I}_{3n_c}
\end{bmatrix},
\end{equation}

which corresponds to

\begin{align}
\mathbf{p}_{b,k+1}
&=
\mathbf{p}_{b,k}
+
\Delta t \mathbf{v}_{b,k},
\\
\mathbf{v}_{b,k+1}
&=
\mathbf{v}_{b,k}.
\end{align}

The control-input matrix is

\begin{equation}
\mathbf{B}
=
\begin{bmatrix}
\frac{1}{2}\Delta t^2 \mathbf{I}_3 \\
\Delta t \mathbf{I}_3 \\
\mathbf{0}
\end{bmatrix}.
\end{equation}

Hence,

\begin{align}
\mathbf{p}_{b,k+1}
&=
\mathbf{p}_{b,k}
+
\Delta t \mathbf{v}_{b,k}
+
\frac{1}{2}\Delta t^2 \mathbf{a}_k,
\\
\mathbf{v}_{b,k+1}
&=
\mathbf{v}_{b,k}
+
\Delta t \mathbf{a}_k.
\end{align}

\subsection{Measurement Model}

The observation vector is constructed as

\begin{equation}
\mathbf{y}
=
\begin{bmatrix}
\mathbf{p}_s \
\mathbf{v}_s \
\mathbf{h}_f
\end{bmatrix},
\end{equation}

where

\begin{itemize}
\item $(\mathbf{p}_s)$ is the foot position measurement,
\item $(\mathbf{v}_s)$ is the foot velocity measurement,
\item $(\mathbf{h}_f)$ is the measured foot height.
\end{itemize}

The code computes

\begin{align}
\mathbf{p}_s &= -\mathbf{p}_{foot}, \\
\mathbf{v}_s &= -\mathbf{v}_{foot}.
\end{align}

The observation model is

\begin{equation}
\mathbf{y}
=
\mathbf{C}\mathbf{x} + \mathbf{v},
\end{equation}

with

\begin{equation}
\mathbf{C} = \begin{bmatrix}
\mathbf{I}_3 & \mathbf{0} & -\mathbf{I} \\
\mathbf{0} & \mathbf{I}_3 & \mathbf{0} \\
\mathbf{0} & \mathbf{0} & \mathbf{H}
\end{bmatrix}.
\end{equation}

and $\mathbf{v} \approx \mathcal{N}(0, \mathbf{R})$ is the measurement noise.

The first block corresponds to

\begin{equation}
\mathbf{p}_{s}
=
\mathbf{p}_b
-
\mathbf{p}_f.
\end{equation}

The second block corresponds to

\begin{equation}
\mathbf{v}_{s}
=
\mathbf{v}_b.
\end{equation}

The final block extracts the vertical foot position

\begin{equation}
h_i
=
z_{f_i}.
\end{equation}

\subsection{Process Noise Covariance}

The process noise covariance is

\begin{equation}
\mathbf{Q}
=
\mathrm{diag}
\left(
Q_p,
Q_v,
Q_f
\right),
\end{equation}

where

\begin{align}
Q_p &= \frac{\Delta t}{20} \sigma_p^2 \mathbf{I}_3,\\
Q_v &= \frac{9.81 \Delta t}{20} \sigma_v^2 \mathbf{I}_3, \\
Q_f &= \Delta t \sigma_f^2 \mathbf{I}_{3n_c}.
\end{align}

When a foot is not in contact, its process noise is increased by a large factor to reduce confidence in the corresponding state estimate.

\subsection{Measurement Noise Covariance}

The measurement covariance matrix is

\begin{equation}
\mathbf{R}
=
\mathrm{diag}
\left(
R_p,
R_v,
R_h
\right),
\end{equation}

where

\begin{align}
R_p &=\sigma_{fp}^2 \mathbf{I}_3, \\
R_v &= \sigma_{fv}^2 \mathbf{I}_3,\\
R_h &= \sigma_{fh}^2 \mathbf{I}_3.
\end{align}

For swing feet, the corresponding measurement covariance entries are inflated by a large factor, effectively removing them from the filter update.

\subsection{Kalman Filter Prediction}

The predicted state estimate is

\begin{equation}
\hat{\mathbf{x}}^{-}_k
=
\mathbf{A}\hat{\mathbf{x}}_{k-1}
+
\mathbf{B}\mathbf{a}_k.
\end{equation}

The predicted covariance is

\begin{equation}
\mathbf{P}^{-}_k
=
\mathbf{A}
\mathbf{P}_{k-1}
\mathbf{A}^{T}
+
\mathbf{Q}.
\end{equation}

\subsection{Kalman Filter Update}

The innovation is

\begin{equation}
\mathbf{e}_k
=
\mathbf{y}_k
-
\mathbf{C}
\hat{\mathbf{x}}^{-}_k.
\end{equation}

The innovation covariance is

\begin{equation}
\mathbf{S}_k
=
\mathbf{C}
\mathbf{P}^{-}_k
\mathbf{C}^{T}
+
\mathbf{R}.
\end{equation}

The Kalman gain is

\begin{equation}
\mathbf{K}_k
=
\mathbf{P}^{-}_k
\mathbf{C}^{T}
\mathbf{S}_k^{-1}.
\end{equation}

The corrected state estimate becomes

\begin{equation}
\hat{\mathbf{x}}_k
=
\hat{\mathbf{x}}^{-}_k
+
\mathbf{K}_k
\mathbf{e}_k.
\end{equation}

Finally, the covariance update is

\begin{equation}
\mathbf{P}_k
=
\left(
\mathbf{I}
-\mathbf{K}_k
\mathbf{C}
\right)
\mathbf{P}^{-}_k.
\end{equation}

\subsection{Estimated Robot State}

The estimated floating-base state is extracted as

\begin{align}
\mathbf{p}_b
&=
\hat{\mathbf{x}}(1:3),
\
\mathbf{v}_b
&=
\hat{\mathbf{x}}(4:6).
\end{align}

These values are passed to

\begin{lstlisting}[language=C++]
updateLinear(
xHat_.segment<3>(0),
xHat_.segment<3>(3));
\end{lstlisting}

to update the robot state representation and generate odometry messages.
